\documentclass[10pt,a4paper]{amsart}
\usepackage{amsmath,amssymb,amsthm,graphicx}
% Retained minimal preamble for the Stacks Project chapter more-morphisms.tex (More on Morphisms).
% Verbatim excerpt of the chapter preamble: Xy-pic setup (native preamble lines 14-18) and the native
% theorem declarations with the mathematical macros (native preamble lines 154-217). Copied mechanically;
% nothing added, removed or redefined. The wrapper class and the amsmath/amssymb/amsthm/graphicx base
% packages are supplied by the assembler, so the class/externaldocument part of the native preamble is omitted.
% For commutative diagrams we use Xy-pic
\usepackage[all]{xy}

% We use 2cell for 2-commutative diagrams.
\xyoption{2cell}
% Theorem environments.
%
\theoremstyle{plain}
\newtheorem{theorem}[subsection]{Theorem}
\newtheorem{proposition}[subsection]{Proposition}
\newtheorem{lemma}[subsection]{Lemma}

\theoremstyle{definition}
\newtheorem{definition}[subsection]{Definition}
\newtheorem{example}[subsection]{Example}
\newtheorem{exercise}[subsection]{Exercise}
\newtheorem{situation}[subsection]{Situation}

\theoremstyle{remark}
\newtheorem{remark}[subsection]{Remark}
\newtheorem{remarks}[subsection]{Remarks}

\numberwithin{equation}{subsection}

% Macros
%
\def\lim{\mathop{\mathrm{lim}}\nolimits}
\def\colim{\mathop{\mathrm{colim}}\nolimits}
\def\Spec{\mathop{\mathrm{Spec}}}
\def\Hom{\mathop{\mathrm{Hom}}\nolimits}
\def\Ext{\mathop{\mathrm{Ext}}\nolimits}
\def\SheafHom{\mathop{\mathcal{H}\!\mathit{om}}\nolimits}
\def\SheafExt{\mathop{\mathcal{E}\!\mathit{xt}}\nolimits}
\def\Sch{\mathit{Sch}}
\def\Mor{\mathop{\mathrm{Mor}}\nolimits}
\def\Ob{\mathop{\mathrm{Ob}}\nolimits}
\def\Sh{\mathop{\mathit{Sh}}\nolimits}
\def\NL{\mathop{N\!L}\nolimits}
\def\CH{\mathop{\mathrm{CH}}\nolimits}
\def\proetale{{pro\text{-}\acute{e}tale}}
\def\etale{{\acute{e}tale}}
\def\QCoh{\mathit{QCoh}}
\def\Ker{\mathop{\mathrm{Ker}}}
\def\Im{\mathop{\mathrm{Im}}}
\def\Coker{\mathop{\mathrm{Coker}}}
\def\Coim{\mathop{\mathrm{Coim}}}

% Boxtimes
%
\DeclareMathSymbol{\boxtimes}{\mathbin}{AMSa}{"02}

%
% Macros for moduli stacks/spaces
%
\def\QCohstack{\mathcal{QC}\!\mathit{oh}}
\def\Cohstack{\mathcal{C}\!\mathit{oh}}
\def\Spacesstack{\mathcal{S}\!\mathit{paces}}
\def\Quotfunctor{\mathrm{Quot}}
\def\Hilbfunctor{\mathrm{Hilb}}
\def\Curvesstack{\mathcal{C}\!\mathit{urves}}
\def\Polarizedstack{\mathcal{P}\!\mathit{olarized}}
\def\Complexesstack{\mathcal{C}\!\mathit{omplexes}}
% \Pic is the operator that assigns to X its picard group, usage \Pic(X)
% \Picardstack_{X/B} denotes the Picard stack of X over B
% \Picardfunctor_{X/B} denotes the Picard functor of X over B
\def\Pic{\mathop{\mathrm{Pic}}\nolimits}
\def\Picardstack{\mathcal{P}\!\mathit{ic}}
\def\Picardfunctor{\mathrm{Pic}}
\def\Deformationcategory{\mathcal{D}\!\mathit{ef}}
\usepackage{cleveref}
\crefname{originalreference}{original reference}{original references}
\Crefname{originalreference}{Original reference}{Original references}
\makeatletter
\newlabel{morphisms-lemma-characterize-affine}{{01S8}{source}{}{sourcepointer-1}{}}
\newlabel{morphisms-lemma-characterize-affine@cref}{{[originalreference][0][]01S8}{[1][0][]source}{}{}{}}
\makeatother
\crefname{originalreference}{original reference}{original references}
\Crefname{originalreference}{Original reference}{Original references}
\makeatletter
\newlabel{properties-lemma-Noetherian-irreducible-components}{{0BA8}{source}{}{sourcepointer-2}{}}
\newlabel{properties-lemma-Noetherian-irreducible-components@cref}{{[originalreference][0][]0BA8}{[1][0][]source}{}{}{}}
\makeatother
\crefname{originalreference}{original reference}{original references}
\Crefname{originalreference}{Original reference}{Original references}
\makeatletter
\newlabel{properties-lemma-locally-Noetherian-dimension-0}{{0AAX}{source}{}{sourcepointer-3}{}}
\newlabel{properties-lemma-locally-Noetherian-dimension-0@cref}{{[originalreference][0][]0AAX}{[1][0][]source}{}{}{}}
\makeatother
\crefname{originalreference}{original reference}{original references}
\Crefname{originalreference}{Original reference}{Original references}
\makeatletter
\newlabel{properties-lemma-point-and-maximal-points-affine}{{01ZX}{source}{}{sourcepointer-4}{}}
\newlabel{properties-lemma-point-and-maximal-points-affine@cref}{{[originalreference][0][]01ZX}{[1][0][]source}{}{}{}}
\makeatother
\title{The affine stratification number of a Noetherian scheme is bounded by its dimension}
\author{The Stacks Project Authors}
\date{stacks-project:more-morphisms.tex:local-snapshot-20261011; GNU FDL 1.2 or later}
\begin{document}
\maketitle
\noindent\textit{Editorial note.} Counted span is the author's own proof environment, source lines 21613-21636, immediately following the statement at 21607-21611.\par\medskip
\noindent\textit{Editorial note.} The reason for dim(Z) < dim(X) (source line 21629 asserts it without proof) is recorded in proof\_roadmap as the reviewer's checklist; the author text is retained verbatim.\par\medskip
\noindent\textit{Editorial note.} All pointers are cross-chapter (Properties x3, Morphisms x1), recorded in assembly.external\_references.\par\medskip
\noindent\textit{Allowed standard prerequisites (editorial).}\begin{itemize}
\item Dimension zero Noetherian schemes are affine. Properties, Lemma properties-lemma-locally-Noetherian-dimension-0, cited at proof lines 21612-21613. Base case d = 0 of the induction.
\item Generic points of irreducible components. Properties, Lemma properties-lemma-Noetherian-irreducible-components, cited at proof lines 21614-21616. Lists the finitely many generic points eta\_1,...,eta\_n that the affine opens must contain.
\item Affine opens through a point and its generalizations. Properties, Lemma properties-lemma-point-and-maximal-points-affine, cited at proof lines 21616-21618 and 21622-21624. Gives affine opens containing all the generic points, used twice (cover step and intersection step).
\item Affineness detected on an affine cover. Morphisms, Lemma morphisms-lemma-characterize-affine, cited at proof lines 21625-21627. U -> X is affine because U -> U\_j is affine for every member of the finite affine cover.
\end{itemize}\medskip
\ifdefined\hypertarget\hypertarget{sourcepointer-1}{}\fi
\noindent\textit{Source reference.} Original reference 01S8: focused-seeds/geometry/S9-morph-asn-bound-noetherian/ref-tags.txt, source lines 5--5. Navigation pointer only: the counted proof cites this statement by \textbackslash{}ref as a standing input; the deduction the counted proof has to make is carried out inside the counted span, and no part of that deduction is delegated to this pointer.\par\medskip
\ifdefined\hypertarget\hypertarget{sourcepointer-2}{}\fi
\noindent\textit{Source reference.} Original reference 0BA8: focused-seeds/geometry/S9-morph-asn-bound-noetherian/ref-tags.txt, source lines 6--6. Navigation pointer only: the counted proof cites this statement by \textbackslash{}ref as a standing input; the deduction the counted proof has to make is carried out inside the counted span, and no part of that deduction is delegated to this pointer.\par\medskip
\ifdefined\hypertarget\hypertarget{sourcepointer-3}{}\fi
\noindent\textit{Source reference.} Original reference 0AAX: focused-seeds/geometry/S9-morph-asn-bound-noetherian/ref-tags.txt, source lines 7--7. Navigation pointer only: the counted proof cites this statement by \textbackslash{}ref as a standing input; the deduction the counted proof has to make is carried out inside the counted span, and no part of that deduction is delegated to this pointer.\par\medskip
\ifdefined\hypertarget\hypertarget{sourcepointer-4}{}\fi
\noindent\textit{Source reference.} Original reference 01ZX: focused-seeds/geometry/S9-morph-asn-bound-noetherian/ref-tags.txt, source lines 8--8. Navigation pointer only: the counted proof cites this statement by \textbackslash{}ref as a standing input; the deduction the counted proof has to make is carried out inside the counted span, and no part of that deduction is delegated to this pointer.\par\medskip
\begin{lemma}
\label{lemma-affine-stratification-number-bound-Noetherian}
Let $X$ be a Noetherian scheme of dimension $\infty > d \geq 0$.
Then the affine stratification number of $X$ is at most $d$.
\end{lemma}

\begin{proof}
By induction on $d$. If $d = 0$, then $X$ is affine, see
Properties, Lemma \ref{properties-lemma-locally-Noetherian-dimension-0}.
Assume $d > 0$. Let $\eta_1, \ldots, \eta_n$ be the generic points
of the irreducible components of $X$ (Properties, Lemma
\ref{properties-lemma-Noetherian-irreducible-components}).
We can cover $X$ by affine opens containing $\eta_1, \ldots, \eta_n$, see
Properties, Lemma \ref{properties-lemma-point-and-maximal-points-affine}.
Since $X$ is quasi-compact we can find a finite affine open covering
$X = \bigcup_{j = 1, \ldots, m} U_j$ with
$\eta_1, \ldots, \eta_n \in U_j$ for all $j = 1, \ldots, m$.
Choose an affine open $U \subset U_1 \cap \ldots \cap U_m$
containing $\eta_1, \ldots, \eta_n$ (possible by the lemma
already quoted). Then the morphism $U \to X$ is affine
because $U \to U_j$ is affine for all $j$, see
Morphisms, Lemma \ref{morphisms-lemma-characterize-affine}.
Let $Z = X \setminus U$.
By construction $\dim(Z) < \dim(X)$.
By induction hypothesis we can find an
affine stratification $Z = \bigcup_{i \in \{0, \ldots, n\}} Z_i$
of $Z$ with $n \leq \dim(Z)$.
Setting $U = X_{n + 1}$ and $X_i = Z_i$ for $i \leq n$
we conclude.
\end{proof}

\end{document}
